\documentclass[11pt,a4paper]{article}
\usepackage[utf8]{inputenc}
\usepackage[T1]{fontenc}
\usepackage[french]{babel}
\usepackage{amsmath, amssymb, amsthm}
\usepackage{geometry}
\usepackage{xcolor}
\geometry{hmargin=2cm, vmargin=2cm}

% Macro pour formater le raisonnement
\newcommand{\raisonnement}[1]{
    \vspace{0.2cm}
    \noindent\textcolor{blue!80!black}{
        \textbf{\textit{Raisonnement : }} \textit{#1}
    }
    \vspace{0.2cm}
}

\begin{document}

\begin{center}
    \textbf{\Large Corrigé de l'Examen MT330} \\
    \vspace{0.2cm}
    \textbf{GRENOBLE INP Esisar UGA, VALENCE}
\end{center}

\vspace{0.5cm}

\section*{Exercice 1 : QCM (5 points)}

\textbf{1. Soient X et Y deux variables aléatoires indépendantes telles que X suit une loi exponentielle de paramètre $\alpha$ et Y suit une loi exponentielle de paramètre $\beta$. $P(X<Y)$ est égale à :}
\textbf{Réponse :} $\frac{\alpha}{\alpha+\beta}$
\raisonnement{C'est une propriété classique des lois exponentielles indépendantes : la probabilité que la première soit inférieure à la seconde est le ratio de son propre paramètre sur la somme des deux paramètres.}

\textbf{2. Énoncé du théorème de transfert :}
\textbf{Réponse :} $\mathbb{E}(\varphi(X)) = \int_{-\infty}^{+\infty} \varphi(t)f(t)dt$
\raisonnement{Le théorème de transfert permet de calculer l'espérance d'une fonction d'une variable aléatoire directement à partir de la densité de cette dernière, sans avoir besoin de déterminer la loi de $\varphi(X)$.}

\textbf{3. Soit X une variable aléatoire qui suit une loi exponentielle de paramètre $\lambda$. Alors, la variance de X est égale à :}
\textbf{Réponse :} $V(X) = \frac{1}{\lambda^2}$
\raisonnement{C'est un résultat de cours : pour une loi exponentielle $\mathcal{E}(\lambda)$, l'espérance est $1/\lambda$ et la variance est $1/\lambda^2$.}

\textbf{4. Soit X une variable aléatoire qui suit une loi normale de moyenne m et d'écart-type $\sigma$, $X \sim \mathcal{N}(m,\sigma)$. La densité de X est :}
\textbf{Réponse :} $f(x) = \frac{1}{\sqrt{2\pi}\sigma} e^{-(x-m)^2/(2\sigma^2)}$
\raisonnement{L'énoncé précise explicitement que le second paramètre de la notation est l'écart-type $\sigma$. La formule classique de la densité de la loi normale comporte $\sigma$ en dehors de l'exponentielle et $2\sigma^2$ au dénominateur dans l'exponentielle.}

\textbf{5. Soient X et Y deux variables indépendantes telles que $X \sim \mathcal{N}(50;3)$ et $Y \sim \mathcal{N}(60;4)$. Alors :}
\textbf{Réponse :} $X+Y \sim \mathcal{N}(110;5)$
\raisonnement{Puisque la notation de l'énoncé précédent est $\mathcal{N}(m, \sigma)$, les écart-types sont 3 et 4. Les variances s'additionnent par indépendance : $V(X+Y) = 3^2 + 4^2 = 9 + 16 = 25$. L'écart-type de la somme est donc $\sqrt{25} = 5$. La moyenne est $50 + 60 = 110$.}

\textbf{6. Soit X une variable aléatoire qui suit une loi exponentielle de paramètre $\lambda$, et soit $Y=X^2$. La densité de Y est :}
\textbf{Réponse :} $f(x) = \frac{\lambda}{2\sqrt{x}} e^{-\lambda\sqrt{x}}$ si $x>0$
\raisonnement{La fonction de répartition de Y est $P(Y \le y) = P(X \le \sqrt{y}) = 1 - e^{-\lambda\sqrt{y}}$. En dérivant cette fonction par rapport à y pour obtenir la densité, on applique la règle de la chaîne, ce qui fait apparaître le facteur $1/(2\sqrt{y})$. (Note : l'énoncé utilise la variable muette $x$ dans les propositions, le principe reste le même).}

\vspace{0.5cm}
\hrule
\vspace{0.5cm}

\section*{Exercice 2 (6 points)}

Soit $f(x) = 6(a-x^2)$ pour $x \in [-\sqrt{a}, \sqrt{a}]$, et 0 sinon.

\textbf{1. Déterminer la valeur de a pour que f soit une densité.}
\raisonnement{L'intégrale de la densité sur son domaine de définition doit être égale à 1.}
\begin{align*}
\int_{-\sqrt{a}}^{\sqrt{a}} 6(a-x^2) dx &= 6 \left[ ax - \frac{x^3}{3} \right]_{-\sqrt{a}}^{\sqrt{a}} \\
&= 6 \left( \left(a\sqrt{a} - \frac{a\sqrt{a}}{3}\right) - \left(-a\sqrt{a} + \frac{a\sqrt{a}}{3}\right) \right) \\
&= 6 \left( \frac{2}{3}a\sqrt{a} - \left(-\frac{2}{3}a\sqrt{a}\right) \right) = 6 \left( \frac{4}{3}a\sqrt{a} \right) = 8a\sqrt{a}
\end{align*}
On résout $8a^{3/2} = 1$, ce qui donne $a^{3/2} = 1/8$, soit $\mathbf{a = 1/4}$.
Le domaine est donc $[-1/2, 1/2]$.

\textbf{2. Calculer l'espérance et la variance de X.}
\raisonnement{La fonction de densité est paire et son intervalle d'intégration est symétrique par rapport à 0, l'espérance est donc nulle.}
$$ \mathbb{E}(X) = \int_{-1/2}^{1/2} x \cdot 6\left(\frac{1}{4}-x^2\right) dx = \mathbf{0} $$

Pour la variance : $V(X) = \mathbb{E}(X^2) - (\mathbb{E}(X))^2$.
\begin{align*}
V(X) &= \int_{-1/2}^{1/2} x^2 \cdot 6\left(\frac{1}{4}-x^2\right) dx = 6 \int_{-1/2}^{1/2} \left(\frac{x^2}{4} - x^4\right) dx \\
&= 6 \left[ \frac{x^3}{12} - \frac{x^5}{5} \right]_{-1/2}^{1/2} = 12 \left( \frac{1/8}{12} - \frac{1/32}{5} \right) \\
&= 12 \left( \frac{1}{96} - \frac{1}{160} \right) = 12 \left( \frac{5 - 3}{480} \right) = \frac{24}{480} = \mathbf{\frac{1}{20}}
\end{align*}

\textbf{3. Soit $Y=X^2$, déterminer la loi de Y.}
\raisonnement{On passe par la fonction de répartition de Y. Puisque $X \in [-1/2, 1/2]$, $Y \in [0, 1/4]$.}
Pour $y \in [0, 1/4]$ :
\begin{align*}
F_Y(y) &= P(X^2 \le y) = P(-\sqrt{y} \le X \le \sqrt{y}) = \int_{-\sqrt{y}}^{\sqrt{y}} 6\left(\frac{1}{4}-x^2\right) dx \\
&= 12 \left[ \frac{x}{4} - \frac{x^3}{3} \right]_0^{\sqrt{y}} = 3\sqrt{y} - 4y\sqrt{y} = 3y^{1/2} - 4y^{3/2}
\end{align*}
La densité de $Y$ s'obtient en dérivant $F_Y(y)$ :
$$ f_Y(y) = \mathbf{\frac{3}{2\sqrt{y}} - 6\sqrt{y}} \quad \text{si } y \in \left]0, \frac{1}{4}\right[ \text{ et } 0 \text{ sinon.} $$

\vspace{0.5cm}
\hrule
\vspace{0.5cm}

\section*{Exercice 3 (6 points)}

Densité jointe : $f(x,y) = k\left(\frac{1}{x^2} + y^2\right)$ pour $x \in [1, 5]$ et $y \in [-1, 1]$.

\textbf{1. Pour quelle valeur de k la fonction f est-elle une densité ?}
\raisonnement{L'intégrale double sur le domaine doit valoir 1.}
\begin{align*}
\iint f(x,y) dx dy &= \int_1^5 \left( \int_{-1}^1 k(x^{-2} + y^2) dy \right) dx = k \int_1^5 \left[ \frac{y}{x^2} + \frac{y^3}{3} \right]_{-1}^1 dx \\
&= k \int_1^5 \left( \frac{2}{x^2} + \frac{2}{3} \right) dx = k \left[ -\frac{2}{x} + \frac{2x}{3} \right]_1^5 \\
&= k \left( \left(-\frac{2}{5} + \frac{10}{3}\right) - \left(-2 + \frac{2}{3}\right) \right) = k \left( \frac{-6+50}{15} - \frac{-6+2}{3} \right) \\
&= k \left( \frac{44}{15} + \frac{4}{3} \right) = k \left( \frac{44}{15} + \frac{20}{15} \right) = k \frac{64}{15}
\end{align*}
On obtient $\mathbf{k = \frac{15}{64}}$.

\textbf{2. Déterminer les lois marginales de X et de Y.}
\raisonnement{On intègre la densité jointe par rapport à l'autre variable.}
Pour $x \in [1, 5]$ :
$$ f_X(x) = \int_{-1}^1 \frac{15}{64} \left( \frac{1}{x^2} + y^2 \right) dy = \mathbf{\frac{15}{32x^2} + \frac{5}{32}} $$
Pour $y \in [-1, 1]$ :
$$ f_Y(y) = \int_1^5 \frac{15}{64} \left( \frac{1}{x^2} + y^2 \right) dx = \frac{15}{64} \left[ -\frac{1}{x} + xy^2 \right]_1^5 = \mathbf{\frac{3}{16} + \frac{15}{16}y^2} $$

\textbf{3. Les variables X et Y sont-elles indépendantes ?}
\raisonnement{Pour qu'elles le soient, il faut que $f(x,y) = f_X(x)f_Y(y)$.}
Le produit des marginales donne des termes croisés (comme du $1/x^2$ et du $y^2$ et des constantes) qui ne se factorisent pas de manière identique à $f(x,y)$.
\textbf{Non}, elles ne sont pas indépendantes.

\textbf{4. Calculer la covariance de X et Y.}
\raisonnement{On sait que $cov(X,Y) = \mathbb{E}(XY) - \mathbb{E}(X)\mathbb{E}(Y)$. Observons les parités par rapport à y.}
La densité marginale $f_Y(y)$ est paire sur $[-1, 1]$, donc $\mathbb{E}(Y) = 0$.
De plus, $\mathbb{E}(XY) = \int_1^5 \int_{-1}^1 xy f(x,y) dy dx$. L'intégrande est impair par rapport à y sur un intervalle symétrique, donc l'intégrale interne vaut 0, d'où $\mathbb{E}(XY) = 0$.
Ainsi, $\mathbf{cov(X,Y) = 0}$.

\vspace{0.5cm}
\hrule
\vspace{0.5cm}

\section*{Exercice 4 (4 points)}

Soient $X, Y, Z \sim \mathcal{N}(0,1)$ indépendantes. $U = X^2 + Y^2 + Z^2$.

\textbf{1. Quelle est la loi de U ? Espérance et variance.}
\raisonnement{C'est une définition classique : la somme des carrés de n lois normales centrées réduites indépendantes suit une loi du Khi-deux à n degrés de liberté.}
U suit une \textbf{loi du $\chi^2$ à 3 degrés de liberté} ($\chi^2_3$).
Ses caractéristiques sont : $\mathbf{\mathbb{E}(U) = 3}$ et $\mathbf{V(U) = 6}$.

\textbf{2. Fonction de répartition et densité de U.}
\raisonnement{On utilise l'indication de l'énoncé sur les coordonnées sphériques. L'événement $(U \le u)$ correspond à la boule de rayon $\sqrt{u}$ centrée à l'origine.}
La densité jointe est $f(x,y,z) = \frac{1}{(2\pi)^{3/2}} e^{-(x^2+y^2+z^2)/2}$.
\begin{align*}
F_U(u) &= \iiint_{x^2+y^2+z^2 \le u} \frac{1}{(2\pi)^{3/2}} e^{-(x^2+y^2+z^2)/2} dx dy dz \\
&= \int_0^{\sqrt{u}} \int_0^{2\pi} \int_0^\pi \frac{1}{(2\pi)^{3/2}} e^{-\rho^2/2} \rho^2 \sin(\varphi) d\varphi d\theta d\rho
\end{align*}
L'intégration angulaire donne $\int_0^{2\pi} d\theta \int_0^\pi \sin(\varphi) d\varphi = 2\pi \times 2 = 4\pi$.
Donc $F_U(u) = \frac{4\pi}{(2\pi)^{3/2}} \int_0^{\sqrt{u}} \rho^2 e^{-\rho^2/2} d\rho = \mathbf{\sqrt{\frac{2}{\pi}} \int_0^{\sqrt{u}} \rho^2 e^{-\rho^2/2} d\rho}$.

En dérivant selon $u$ avec le théorème fondamental de l'analyse et la règle de composition (dérivée de la borne $\sqrt{u}$ qui est $1/(2\sqrt{u})$) :
$$ f_U(u) = F_U'(u) = \sqrt{\frac{2}{\pi}} (\sqrt{u})^2 e^{-(\sqrt{u})^2/2} \cdot \frac{1}{2\sqrt{u}} = \mathbf{\frac{1}{\sqrt{2\pi}} \sqrt{u} e^{-u/2}} \quad \text{pour } u>0. $$

\vspace{0.5cm}
\hrule
\vspace{0.5cm}

\section*{Exercice 5 (3 points)}

Densité $f(x) = \frac{1}{2\theta\sqrt{x}} e^{-\sqrt{x}/\theta}$ pour $x>0$.
Estimation de $\theta$ par le maximum de vraisemblance.

\raisonnement{On écrit la fonction de vraisemblance (produit des densités), on passe au logarithme pour simplifier le produit en somme, on dérive par rapport à $\theta$ et on cherche le zéro de cette dérivée.}
Fonction de vraisemblance :
$$ L(\theta) = \prod_{i=1}^n \frac{1}{2\theta\sqrt{x_i}} e^{-\sqrt{x_i}/\theta} = \frac{1}{(2\theta)^n \prod_{i=1}^n \sqrt{x_i}} \exp\left(-\frac{1}{\theta}\sum_{i=1}^n \sqrt{x_i}\right) $$
Log-vraisemblance :
$$ \ln L(\theta) = -n\ln(2) - n\ln(\theta) - \sum_{i=1}^n \ln(\sqrt{x_i}) - \frac{1}{\theta} \sum_{i=1}^n \sqrt{x_i} $$
Dérivée par rapport à $\theta$ :
$$ \frac{\partial \ln L}{\partial \theta} = -\frac{n}{\theta} + \frac{1}{\theta^2} \sum_{i=1}^n \sqrt{x_i} $$
On annule la dérivée :
$$ -\frac{n}{\theta} + \frac{1}{\theta^2} \sum_{i=1}^n \sqrt{x_i} = 0 \iff \frac{n}{\theta} = \frac{1}{\theta^2} \sum_{i=1}^n \sqrt{x_i} \iff n\theta = \sum_{i=1}^n \sqrt{x_i} $$
L'estimateur est donc : $\mathbf{\hat{\theta} = \frac{1}{n} \sum_{i=1}^n \sqrt{x_i}}$.

\vspace{0.5cm}
\hrule
\vspace{0.5cm}

\section*{Exercice 6 (8 points)}

\textbf{1(a). Calculer la probabilité de l'événement « Robert tombe sur la piste qu'il a choisie ».}
\raisonnement{On applique la formule des probabilités totales. Soient B, R, N les choix des pistes (Bleue, Rouge, Noire) et T l'événement "tomber".}
\begin{align*}
P(T) &= P(T|B)P(B) + P(T|R)P(R) + P(T|N)P(N) \\
&= \left(\frac{1}{10} \times \frac{1}{4}\right) + \left(\frac{1}{6} \times \frac{1}{2}\right) + \left(\frac{2}{5} \times \frac{1}{4}\right) \\
&= \frac{1}{40} + \frac{1}{12} + \frac{1}{10} = \frac{3 + 10 + 12}{120} = \mathbf{\frac{5}{24}}
\end{align*}

\textbf{1(b). Quelle est la probabilité que Robert ait choisi la piste noire sachant qu'il est tombé ?}
\raisonnement{C'est une probabilité conditionnelle inverse (théorème de Bayes).}
$$ P(N|T) = \frac{P(T|N)P(N)}{P(T)} = \frac{\frac{2}{5} \times \frac{1}{4}}{\frac{5}{24}} = \frac{1/10}{5/24} = \frac{24}{50} = \mathbf{0.48} $$

\textbf{2(a). Loi de $Z = X+Y$.}
\raisonnement{On calcule le produit de convolution entre $X \sim \mathcal{E}(1/5)$ et $Y \sim \mathcal{U}[10; 15]$.}
$f_Z(z) = \int_{10}^{15} f_X(z-y)f_Y(y) dy$. Trois cas se présentent pour le support de X :
\begin{itemize}
    \item Si $z < 10$ : $\mathbf{f_Z(z) = 0}$.
    \item Si $10 \le z \le 15$ : on intègre de 10 à z. $f_Z(z) = \frac{1}{25} e^{-z/5} \int_{10}^z e^{y/5} dy = \mathbf{\frac{1}{5} \left( 1 - e^{-(z-10)/5} \right)}$.
    \item Si $z > 15$ : on intègre de 10 à 15. $f_Z(z) = \frac{1}{25} e^{-z/5} \int_{10}^{15} e^{y/5} dy = \mathbf{\frac{1}{5} e^{-z/5} \left( e^3 - e^2 \right)}$.
\end{itemize}

\textbf{2(b). Probabilité qu'il ne soit pas en retard.}
\raisonnement{Il a 15 minutes devant lui. L'événement est $Z \le 15$.}
\begin{align*}
P(Z \le 15) &= \int_{10}^{15} \frac{1}{5} \left( 1 - e^{-(z-10)/5} \right) dz = \left[ \frac{z}{5} + e^{-(z-10)/5} \right]_{10}^{15} \\
&= \left( 3 + e^{-1} \right) - (2 + 1) = \mathbf{e^{-1}} \approx 0.368
\end{align*}

\textbf{2(c). Probabilité que Stéphanie arrive avant Robert.}
\raisonnement{Les variables sont $T_1 \sim \mathcal{U}[55; 85]$ (Robert) et $T_2 \sim \mathcal{U}[65; 75]$ (Stéphanie). On cherche $P(T_2 < T_1)$.}
L'aire totale du domaine (rectangle) est $30 \times 10 = 300$.
Le sous-domaine favorable correspond à $t_2 < t_1$. Ce domaine se découpe en deux :
- Un triangle pour $t_1 \in [65; 75]$, dont l'aire est $\frac{10 \times 10}{2} = 50$.
- Un rectangle plein pour $t_1 \in [75; 85]$, d'aire $10 \times 10 = 100$.
L'aire favorable totale est $150$. La probabilité est donc $\frac{150}{300} = \mathbf{\frac{1}{2}}$.

\textbf{3(a). Loi exacte de $S_n$.}
\raisonnement{On répète n expériences indépendantes avec deux issues (succès = skis Castafiore).}
$S_n$ suit une \textbf{loi binomiale} de paramètres $n$ et $p$ : $\mathcal{B}(n, p)$.

\textbf{3(b). Valeur de n pour déterminer p à $10^{-2}$ près avec risque $0.05$.}
\raisonnement{La demi-largeur de l'intervalle de confiance est fixée à 0.01. Sans information sur p, on se place dans le cas défavorable où $p=0.5$ maximisant la variance.}
$$ 1.96 \sqrt{\frac{0.5 \times 0.5}{n}} \le 0.01 \iff \sqrt{n} \ge \frac{1.96 \times 0.5}{0.01} = 98 $$
Soit $n \ge 98^2 = \mathbf{9604}$ skieurs.

\textbf{3(c). Intervalle de confiance avec 1200 skieurs et 250 skis Castafiore.}
\raisonnement{On applique la formule classique de l'intervalle de confiance pour une proportion.}
Fréquence estimée : $\hat{p} = \frac{250}{1200} = \frac{5}{24} \approx 0.208$.
La marge d'erreur $E$ vaut $1.96 \sqrt{\frac{\hat{p}(1-\hat{p})}{n}} = 1.96 \sqrt{\frac{(5/24) \times (19/24)}{1200}} \approx 1.96 \times 0.0117 \approx 0.023$.
L'intervalle de confiance au risque 5\% est donc $\approx [0.208 - 0.023, 0.208 + 0.023] = \mathbf{[0.185 ; 0.231]}$.

\end{document}